% =====================================================================
%  Pearls AQI Predictor - Project Report
%  10Pearls Shine Internship, Data Science track
%
%  Compile with:  pdflatex PROJECT_REPORT.tex   (run twice for the ToC)
% =====================================================================
\documentclass[11pt,a4paper]{article}

\usepackage[margin=2.4cm]{geometry}
\usepackage{booktabs}
\usepackage{graphicx}
\usepackage{xcolor}
\usepackage{titlesec}
\usepackage{fancyhdr}
\usepackage{enumitem}
\usepackage{amssymb}
\usepackage{needspace}
\usepackage{etoolbox}
\usepackage{caption}
\usepackage[hidelinks]{hyperref}
\hypersetup{
  pdftitle={Pearls AQI Predictor: a three-day air quality forecasting system for Ghotki, Sindh},
  pdfauthor={Shivam Kumar},
  pdfsubject={10Pearls Shine Internship, Data Science Track},
  pdfkeywords={air quality, AQI, forecasting, machine learning, Ghotki}
}

\definecolor{ink}{HTML}{1A1A1A}
\definecolor{accent}{HTML}{6FA82B}
\definecolor{muted}{HTML}{5F5F5A}
\definecolor{rule}{HTML}{D8D8D4}

\titlespacing*{\section}{0pt}{16pt}{7pt}
\titlespacing*{\subsection}{0pt}{12pt}{5pt}

% Keep a heading with at least a few lines of what follows it, so no
% section title is stranded at the foot of a page. The reservation goes
% in titlesec's own before-code: a \pretocmd hook on \section does not
% fire, because titlesec has already redefined the command.
\titleformat{\section}[block]
  {\Needspace*{10\baselineskip}\Large\bfseries\color{ink}}
  {\thesection}{0.6em}{}
\titleformat{\subsection}[block]
  {\Needspace*{8\baselineskip}\large\bfseries\color{ink}}
  {\thesubsection}{0.6em}{}

\pagestyle{fancy}
\fancyhf{}
\renewcommand{\headrulewidth}{0.4pt}
\lhead{\small\color{muted}Pearls AQI Predictor}
\rhead{\small\color{muted}Ghotki, Sindh}
\cfoot{\small\color{muted}\thepage}

\widowpenalty=10000
\clubpenalty=10000
\displaywidowpenalty=10000
\raggedbottom

\captionsetup{font=small,labelfont={bf,color=ink},skip=6pt}
\setlist[itemize]{leftmargin=1.3em,itemsep=2pt,topsep=4pt}
\setlength{\parskip}{6pt}
\setlength{\parindent}{0pt}

\newcommand{\keyfig}[2]{{\color{accent}\textbf{#1}}~#2}

\begin{document}

% ---------------------------------------------------------------
\begin{titlepage}
\centering
\vspace*{3.2cm}

{\Huge\bfseries\color{ink} Pearls AQI Predictor\par}
\vspace{0.5cm}
{\Large\color{muted} A three-day air quality forecasting system\par}
{\Large\color{muted} for Ghotki, Sindh\par}

\vspace{1.6cm}
{\color{rule}\rule{0.45\textwidth}{0.8pt}}
\vspace{1.6cm}

{\large\bfseries\color{ink} Shivam Kumar\par}
\vspace{0.25cm}
{\normalsize\color{muted} Institute of Business Administration, Karachi\par}

\vspace{1.1cm}

{\large 10Pearls Shine Internship\par}
{\large Data Science Track\par}

\vspace{2cm}

\begin{tabular}{rl}
\textbf{Dashboard} & \href{https://10pearls-aqi-predictor-52wav2lhgysbpemqnqu63e.streamlit.app/}{\small 10pearls-aqi-predictor-52wav2lhgysbpemqnqu63e.streamlit.app} \\[5pt]
\textbf{Repository} & \href{https://github.com/Shivam-kumar-30012/10Pearls-AQI-Predictor}{\small github.com/Shivam-kumar-30012/10Pearls-AQI-Predictor} \\[5pt]
\textbf{Date} & \small September 2026 \\
\end{tabular}

\vfill
\end{titlepage}


% =====================================================================
\section*{Abstract}
\addcontentsline{toc}{section}{Abstract}

This report documents an end-to-end machine learning system that
forecasts the Air Quality Index for Ghotki, Sindh, for every hour of the
next three days. The system collects data hourly, retrains nightly, and
serves predictions through a deployed dashboard, all without manual
intervention.

Three models cover the forecast horizon --- one per day, each predicting
all 24 hours of its block. On a held-out chronological test period they
reach R\textsuperscript{2} of 0.932, 0.716 and 0.613 at one, two and
three days ahead, beating a persistence baseline by a margin that widens
with distance.

Particular attention was given to the target variable. The AQI is
computed to the EPA standard --- correct averaging windows, truncation
before band lookup, and extrapolation rather than clamping above scale
--- and validated independently of model performance, since the target
bounds what any model can achieve. Section~4 sets out the method and
that validation in full.

\subsection*{What this report covers}

\begin{itemize}
  \item \textbf{Sections 1--3} set out the objective, the system
        architecture, and how two independent data sources were combined
        without misalignment.
  \item \textbf{Section 4} sets out how the AQI target is computed, why
        each step of the EPA method is mandatory, and how the
        implementation was validated.
  \item \textbf{Sections 5--6} describe the 71 engineered features and
        the exploratory analysis that justified them.
  \item \textbf{Section 7} covers modelling: problem formulation,
        leakage prevention, the persistence baseline, algorithm
        comparison, three experiments that did not work, and the final
        results.
  \item \textbf{Sections 8--10} cover explainability, the automated
        pipelines and their observed failure modes, and the dashboard.
  \item \textbf{Sections 11--12} state the limitations plainly and draw
        the conclusion.
\end{itemize}

A reader wanting only the outcome should read this abstract and
Section~7.7. A reader interested in the engineering should begin at
Section~4.

\newpage

% =====================================================================
% parskip is suppressed here so the contents fit one page; with the
% document-wide 6pt it spilled a single line onto the next.
{\setlength{\parskip}{0pt}\tableofcontents}
\newpage

% =====================================================================
\section{Objective and scope}

The brief specified a serverless machine learning pipeline with five
stages: fetch raw data from a public API, engineer features into a
feature store, train and evaluate models, register the best of them,
and serve predictions through a web application --- with the feature
pipeline running hourly and training running daily.

Ghotki was selected as the target city. It is a district in northern
Sindh with no dedicated monitoring station of its own, which makes a
model-based forecast more useful there than in a city with dense
instrumentation.

Everything specified was built. Beyond the core requirement, the system
also includes SHAP explainability, hazardous-AQI alerting, and a
documented set of negative results --- experiments that were run
properly and did not work, reported here rather than omitted.

% ---------------------------------------------------------------
\section{System architecture}

\begin{center}
\small
\begin{tabular}{rl}
\toprule
\textbf{Stage} & \textbf{Implementation} \\
\midrule
Raw collection & OpenWeather Air Pollution API, Open-Meteo ERA5 \\
Feature engineering & 71 engineered features, EPA-compliant AQI \\
Feature store & Hopsworks, feature group \texttt{aqi\_ghotki\_features} v3 \\
Training & Ridge, Random Forest, XGBoost --- compared per horizon \\
Registry & Hopsworks Model Registry, currently version 9 \\
Automation & GitHub Actions: hourly collection, nightly training \\
Serving & Streamlit Community Cloud \\
\bottomrule
\end{tabular}
\end{center}

The pipeline is a sequence of numbered scripts, each doing one thing and
writing its output to disk before the next reads it:

\begin{center}
\small
\texttt{01\_backfill\_pollutants} $\rightarrow$
\texttt{02\_backfill\_weather} $\rightarrow$
\texttt{03\_merge\_raw} $\rightarrow$
\texttt{04\_build\_features} \\[3pt]
$\rightarrow$ \texttt{05\_upload\_to\_store} $\rightarrow$
\texttt{06\_train} / \texttt{06b\_train\_delta} $\rightarrow$
\texttt{07\_register\_models} \\[3pt]
$\rightarrow$ \texttt{08\_hourly\_collector} $\rightarrow$
\texttt{09\_predict} $\rightarrow$
\texttt{10\_build\_snapshot} $\rightarrow$
\texttt{11\_explain}
\end{center}

Separating collection from transformation was a deliberate early choice.
Downloading five years of hourly data takes several minutes and is
rate-limited; feature engineering takes seconds. Keeping them apart
meant that every subsequent fix to the feature logic --- including the
one described in Section~\ref{sec:aqi} --- cost thirty seconds rather
than an hour.

% ---------------------------------------------------------------
\section{Data}

\subsection{Sources}

Two APIs, joined on an hourly UTC timestamp.

\textbf{OpenWeather Air Pollution History} supplies PM\textsubscript{2.5},
PM\textsubscript{10}, NO\textsubscript{2}, O\textsubscript{3}, CO,
SO\textsubscript{2} and NH\textsubscript{3} from 27~November~2020, the
earliest date its archive covers for this location.

\textbf{Open-Meteo ERA5} supplies temperature, humidity, pressure, wind
speed, wind direction, cloud cover and precipitation. It was chosen over
OpenWeather's weather endpoint because the free tier of the latter
reaches back only about five days, which is useless for building a
training set. ERA5 is a reanalysis product --- a physics model run
backwards over historical observations --- and is the standard source
for historical weather in research contexts.

\subsection{Combining two sources}

Mixing sources introduces alignment risk, so three things were enforced.
Both feeds are requested in UTC and floored to the hour, removing any
possibility of a timezone offset silently misaligning the join. The join
is an inner join, so a row exists only where both pollutant and weather
data are present --- a left join would have left weather columns empty
on unmatched rows, which then propagates silently into every derived
feature. And OpenWeather's \texttt{-9999} sensor-failure sentinel is
converted to a missing value at fetch time rather than being allowed to
reach a calculation.

The merge yielded \textbf{49,017} hours from 50,016 weather readings and
49,017 pollutant readings. ERA5 has no gaps because it is a model;
OpenWeather has real outages.

\subsection{Gaps and the hourly grid}

The merged data was missing 984 hours (1.97\%), spread across 31 gaps,
the largest lasting 121~hours.

This matters more than the percentage suggests. Lag features are built
with a row shift, so \texttt{shift(24)} means ``24 rows back'', not
``24 hours back''. Across a 121-hour gap, a feature named
\texttt{aqi\_lag\_24h} would silently contain a value from five days
earlier.

The fix is to reindex onto a complete hourly grid before building any
feature, so every row is exactly one hour apart. Gaps of three hours or
less are filled by time-weighted interpolation; longer ones are left
missing and their rows drop out later. The threshold is a judgement:
estimating two missing hours between known values is reasonable,
inventing five days of weather is not.

A detail worth noting is that interpolation \emph{saves} rows rather
than merely patching them. A missing hour does not cost only that hour
--- it breaks the lag features of every row for the following week. Two
interpolated values can preserve around 170 usable rows.

% ---------------------------------------------------------------
\section{Computing the AQI: method and validation}
\label{sec:aqi}

The AQI is the target variable, so its correctness determines the
ceiling on everything downstream. This section sets out how it is
computed, why each step is mandatory, and the validation that confirmed
the implementation was right.

\subsection{The EPA method}

The index is not a single formula. Each pollutant has its own breakpoint
table mapping concentration to an index band; the concentration is
looked up, interpolated linearly within its band, and the overall AQI is
the \emph{maximum} across pollutants --- the dominant-pollutant rule,
which reflects that air is only as good as its worst constituent.

Three requirements in the standard are easy to overlook, and each was
found to matter materially.

\subsubsection*{1. Concentrations must be averaged first}

EPA breakpoints are defined for \emph{averaged} concentrations, not
instantaneous readings:

\begin{center}
\small
\begin{tabular}{ll}
\toprule
\textbf{Pollutant} & \textbf{Averaging window} \\
\midrule
PM\textsubscript{2.5}, PM\textsubscript{10} & 24 hours \\
CO, O\textsubscript{3} & 8 hours \\
\bottomrule
\end{tabular}
\end{center}

The windows are not arbitrary. They match the exposure period each
pollutant's health standard is built on --- particulate harm accrues
over a day, ozone and CO act within hours.

The first implementation passed raw hourly readings into these tables.
The reasoning at the time was that hourly resolution was the point of an
hourly forecast, and averaging appeared to discard information. That
reasoning does not hold, for a specific and instructive reason. Hourly
concentrations are far more volatile than their 24-hour means, and
because the overall AQI takes the maximum across pollutants, that
volatility caused the dominant pollutant to flip between
PM\textsubscript{2.5}, PM\textsubscript{10} and O\textsubscript{3} from
one hour to the next. Consecutive target values were therefore drawn
from unrelated measurements, and the series behaved as though the air
itself were erratic.

\subsubsection*{2. Concentrations must be truncated before lookup}

EPA publishes bands as $(0.0, 12.0)$, $(12.1, 35.4)$, and so on. Read
literally these appear to leave gaps, and a direct implementation of the
published tables will find values that match no band at all.

They are not gaps. The standard requires the concentration to be
\emph{truncated} --- not rounded --- to a fixed number of decimals
before lookup: one decimal for PM\textsubscript{2.5} and CO, zero for
PM\textsubscript{10} and ozone in ppb. After truncation, $12.05$ becomes
$12.0$ and falls cleanly inside the first band. The apparent gaps close
entirely.

Truncation rather than rounding matters: $\mathrm{round}(12.05, 1)$
gives $12.1$, which lands in the \emph{next} band and returns a
different index.

\subsubsection*{3. Above-scale readings should extrapolate, not clamp}

EPA's tables end at defined maxima --- 604\,\textmu g/m\textsuperscript{3}
for PM\textsubscript{10}, for instance. Ghotki experiences dust events
that exceed this; 158 such readings exist in the data.

Clamping them all to 500 is the obvious handling, but it destroys the
gradient: a severe dust event and a catastrophic one become the same
number, and the model loses any ability to distinguish them. The
implementation instead extrapolates along the final band's slope, which
preserves the ordering. Values above 500 are retained for training and
clipped only for display.

\subsection{Validating the implementation}

Because the target bounds every downstream result, the calculator was
validated directly rather than inferred from model performance. The two
are easily conflated, and a model can only ever be as good as its label
allows.

Two checks were run. The first evaluates the calculator at band
boundaries, where any departure from the standard surfaces first:

\begin{center}
\small
\begin{tabular}{lcc}
\toprule
\textbf{Input} & \textbf{Truncated to} & \textbf{AQI} \\
\midrule
PM\textsubscript{2.5} = 12.05 & 12.0 & 50 \\
PM\textsubscript{2.5} = 35.45 & 35.4 & 100 \\
O\textsubscript{3} = 107\,\textmu g/m\textsuperscript{3} & 54\,ppb & 50 \\
PM\textsubscript{10} = 1{,}262 & 1{,}262 & 1{,}158 (extrapolated) \\
Typical Ghotki reading & --- & 99 (PM\textsubscript{2.5}) \\
\bottomrule
\end{tabular}
\end{center}

These cases are asserted as tests in the source, so the behaviour is
locked in rather than merely observed once.

The second check examines the statistical behaviour of the resulting
series against what the physical process implies. Air quality changes
gradually, so a correctly computed index should be smooth and strongly
autocorrelated:

\begin{center}
\begin{tabular}{lr}
\toprule
\textbf{Property of the computed series} & \textbf{Value} \\
\midrule
Median hourly change & 1 \\
Mean hourly change & 1.5 \\
Hourly changes above 100 & 0 \\
Correlation, today vs.\ 24 hours later & 0.797 \\
Dominant pollutant, PM\textsubscript{2.5} & 89.7\% of hours \\
\bottomrule
\end{tabular}
\end{center}

All five behave as the physics predicts. The correlation figure is
particularly useful: squared, it approximates the best
R\textsuperscript{2} any model could achieve against this target, which
sets a realistic expectation of \keyfig{0.66} at 24 hours before any
model is trained. Section~7 reports performance against that
expectation.

The dominant-pollutant split is a further sanity check. PM\textsubscript{2.5}
leading 89.7\% of hours, PM\textsubscript{10} 7.1\% during dust events
and ozone 3.1\% is a physically sensible profile for rural Sindh; a
result distributed evenly across pollutants would indicate the averaging
windows had not been applied.

\subsection{Practices adopted as a result}

Two practices were adopted for the remainder of the project.

\textbf{Regression tests on the calculator.} The specific boundary
values in the table above are asserted in
\texttt{src/aqi\_calculator.py} and run whenever the file is executed
directly. Any future edit that reintroduces the direct-lookup behaviour
fails immediately and visibly.

\textbf{Baseline comparison on every result.} Because the target itself
had proven capable of setting an invisible ceiling, no subsequent score
was accepted without a naive comparison beside it. That is why the
persistence baseline appears in every results table in Section~7, rather
than only in a summary.

\Needspace*{22\baselineskip}
\section{Feature engineering}

The feature table holds \textbf{48,710} rows and \textbf{71} features.

\begin{center}
\small
\begin{tabular}{lll}
\toprule
\textbf{Group} & \textbf{Count} & \textbf{Examples} \\
\midrule
Time and season & 13 & hour, month, four season indicators \\
Cyclical & 4 & \texttt{hour\_sin}, \texttt{month\_cos} \\
AQI lags & 8 & 1\,h through 168\,h \\
AQI rolling & 9 & mean, std, max, min over 1\,d / 3\,d / 7\,d \\
AQI trends & 4 & 6\,h, 24\,h, 72\,h deltas; vs.\ weekly mean \\
Pollutants & 15 & raw, 24\,h lag, 24\,h rolling mean \\
Weather & 7 & temperature, humidity, pressure, wind \\
Weather derived & 11 & wind $u$/$v$, ventilation, inversion proxy \\
\bottomrule
\end{tabular}
\end{center}

Three of these deserve explanation.

\textbf{Wind decomposition.} Wind direction in degrees is circular:
359\textdegree{} and 1\textdegree{} are nearly the same direction but
maximally distant as numbers. Decomposing into east--west and
north--south components fixes this and encodes speed simultaneously. The
wind rose in the exploratory analysis showed why it matters --- air
arriving from the north carries a mean AQI near 160, from the south near
110.

\textbf{Inversion proxy.} \texttt{temp\_range\_24h}, the spread between
the day's highest and lowest temperature, is a proxy for a temperature
inversion. A large swing implies clear skies and strong nocturnal
cooling, which produces a cold layer near the ground that traps smoke
beneath it. It correlates with AQI at $r = +0.30$, and SHAP later
confirmed the day~+1 model uses it heavily.

\textbf{Ventilation index.} Wind speed weighted by dryness, on the
reasoning that wind disperses pollution and dry air does not grow
particles. It correlates at only $r = -0.17$ --- real but modest, and
reported here as such rather than overstated.

Cyclical hour encoding was retained but is likely inert for this target.
Because the AQI is computed from 24-hour rolling averages, hour-of-day
variation is averaged out by construction, and the exploratory analysis
confirms a flat hourly profile. It is noted rather than defended.

% ---------------------------------------------------------------
\section{Exploratory analysis}
\label{sec:eda}

Eight figures were produced. The findings that shaped the modelling:

\textbf{Strong seasonality.} Winter averages near AQI 200, monsoon near
90, across five clean annual cycles. This justifies the season
indicators.

\textbf{PM\textsubscript{2.5} dominates.} It is the deciding pollutant
for 89.7\% of hours, PM\textsubscript{10} for 7.1\% (dust events),
ozone for 3.1\%. A physically sensible profile for rural Sindh.

\textbf{Pressure is the strongest weather signal} at $r = +0.45$. High
pressure implies stagnant air and trapped pollution. Temperature follows
at $r = -0.43$.

\textbf{Bimodal AQI distribution}, with peaks near 80 and 155 ---
consistent with two regimes, clean season and dirty season.

\textbf{Predictability decays smoothly.} Autocorrelation runs 0.997 at
one hour, 0.797 at 24, 0.646 at 48, 0.570 at 72. Computed before any
model was trained, this set the expectation that results were later
compared against, rather than being reverse-justified afterwards.

Time spent in each EPA category: Good 3.3\%, Moderate 34.0\%, Unhealthy
for Sensitive Groups 22.6\%, Unhealthy 30.1\%, Very Unhealthy 8.3\%,
Hazardous 1.6\%.

% ---------------------------------------------------------------
\section{Modelling}

\subsection{Problem formulation}

The requirement is an hourly forecast across three days --- 72 values.
Training 72 separate models would be unwieldy to register and maintain,
so instead three models each cover a 24-hour block, with
\texttt{hours\_ahead} supplied as a feature.

Each source row is expanded into 24 training examples, one per hour in
the block, producing roughly 1.1~million examples per model. The model
therefore learns explicitly that a one-hour forecast is near-certain
while a 24-hour forecast is not.

\subsection{Preventing leakage}

The chronological split is applied to the \emph{source row}, not to the
expanded examples. Splitting the expanded set would allow 3\,pm Tuesday
to appear in training while 4\,pm Tuesday appears in test --- nearly
identical inputs on both sides, producing an inflated score. Splitting
by source row keeps all 24 hours of a day together.

Data is never shuffled. The oldest 80\% trains, the newest 20\% tests,
with the boundary at 2~July~2025. Shuffling would let the model learn
from 2026 to predict 2021.

\subsection{The persistence baseline}

Every result is reported against persistence --- predicting that the AQI
stays exactly as it is now. This matters because AQI is strongly
autocorrelated, so persistence is a demanding benchmark rather than a
trivial one, and a model can post a high R\textsuperscript{2} while
contributing nothing.

\begin{center}
\begin{tabular}{lccc}
\toprule
\textbf{Horizon} & \textbf{Persistence} & \textbf{Model} & \textbf{Gain} \\
\midrule
Day~+1 & 0.908 & 0.932 & $+0.024$ \\
Day~+2 & 0.682 & 0.716 & $+0.034$ \\
Day~+3 & 0.544 & 0.613 & $+0.069$ \\
\bottomrule
\end{tabular}
\end{center}

The gain column is the honest measure of what the models contribute, and
its trend is the strongest result in the project: the advantage
\emph{triples} as the horizon extends. At day~+3 persistence has largely
collapsed while the models hold, which indicates they have learned
genuine atmospheric behaviour rather than simply copying the current
value.

\Needspace*{26\baselineskip}
\subsection{Algorithm comparison}

Ridge, Random Forest and XGBoost were compared at every horizon on
identical data.

\begin{center}
\small
\begin{tabular}{llccc}
\toprule
\textbf{Horizon} & \textbf{Model} & \textbf{R\textsuperscript{2}} & \textbf{RMSE} & \textbf{MAE} \\
\midrule
Day~+1 & \textbf{XGBoost (delta)} & \textbf{0.932} & \textbf{16.5} & \textbf{8.3} \\
       & Ridge                    & 0.930 & 16.8 & 10.2 \\
       & Persistence              & 0.908 & 19.2 & 9.9 \\
       & XGBoost (absolute)       & 0.889 & 21.2 & 9.5 \\
       & Random Forest            & 0.882 & 21.8 & 9.8 \\
\midrule
Day~+2 & \textbf{Ridge}           & \textbf{0.716} & \textbf{33.9} & 22.6 \\
       & Persistence              & 0.682 & 35.8 & 21.5 \\
       & XGBoost                  & 0.652 & 37.5 & 21.5 \\
       & Random Forest            & 0.630 & 38.6 & 22.3 \\
\midrule
Day~+3 & \textbf{Ridge}           & \textbf{0.613} & \textbf{39.5} & 26.1 \\
       & XGBoost                  & 0.595 & 40.4 & 26.2 \\
       & Persistence              & 0.544 & 42.9 & 26.0 \\
       & Random Forest            & 0.529 & 43.5 & 28.4 \\
\bottomrule
\end{tabular}
\end{center}

Ridge winning at the longer horizons is a substantive finding, not a
disappointment. Once the AQI target is computed correctly the series is
smooth and close to linear in its lag features, leaving little curvature
for tree ensembles to exploit. On a noisier target the same outcome
would suggest a weak signal; here it indicates that the problem is
genuinely linear.

\subsection{Experiments that did not work}

Three were run properly and are reported because negative results are
part of the record.

\textbf{Delta targets at days~2 and~3.} Predicting the change rather
than the level helped day~+1 substantially --- MAE fell from 10.2 to 8.3
--- but Ridge scored \emph{identically} at days~2 and~3
(0.7158 vs.\ 0.7158). This is mathematically expected: for a linear
model, predicting $y - x$ when $x$ is already a feature is the same
computation as predicting $y$. XGBoost on the delta target was markedly
worse at those horizons (0.606 and 0.484).

\textbf{Forecast weather instead of observed weather.} A methodologically
important idea: predicting Saturday's AQI from Thursday's weather is
wrong, because Saturday's air depends on Saturday's conditions. Real
forecasting systems train on archived forecasts so that inputs at
training time match what is available at prediction time. Open-Meteo's
previous-runs archive was retrieved, but returned only \textbf{55.5\%}
coverage for wind, pressure and humidity at this location --- precisely
the variables the exploratory analysis identified as most predictive.
Adopting it would have cost more than half the training data in exchange
for incomplete inputs, so observed weather was retained.

\textbf{Interaction terms with lead time.} Adding
\texttt{aqi}~$\times$~\texttt{hours\_ahead} and similar products left
accuracy essentially unchanged (0.9304 $\rightarrow$ 0.9296 at day~+1)
but fixed a real defect: without them the forecast was flat, predicting
an identical value for all 24 hours of a block. They were retained for
that reason rather than for any gain in score.

\subsection{Final results}

\begin{center}
\begin{tabular}{lcccc}
\toprule
\textbf{Horizon} & \textbf{Model} & \textbf{R\textsuperscript{2}} & \textbf{RMSE} & \textbf{MAE} \\
\midrule
Day~+1 & XGBoost (delta) & 0.932 & 16.5 & 8.3 \\
Day~+2 & Ridge           & 0.716 & 33.9 & 22.6 \\
Day~+3 & Ridge           & 0.613 & 39.5 & 26.1 \\
\bottomrule
\end{tabular}
\end{center}

Day-ahead performance is strong: an R\textsuperscript{2} of 0.932 and a
mean absolute error of 8.3 AQI points, against a series whose own mean
is 131. In practical terms the forecast for tomorrow is typically within
a tenth of the scale, and comfortably inside the correct EPA category
for most hours.

Accuracy declines with lead time, as it must. By day~+3 the mean
absolute error has risen to 26 points, which is wide enough to straddle
a category boundary. The daily average remains considerably more
reliable than any individual hour at that distance, which is why the
dashboard leads with the daily figure and treats the hourly curve as
supporting detail.

The decline is not evidence of a modelling shortfall. Autocorrelation in
the series falls from 0.797 at 24 hours to 0.570 at 72
(Section~\ref{sec:eda}), which bounds what is achievable at each
horizon; and persistence --- the strongest naive benchmark for a slow
moving series --- itself only reaches RMSE 42.9 at three days. The
models beat it by their widest margin precisely there.

% ---------------------------------------------------------------
\section{Explainability}

SHAP was applied to all three horizons. It decomposes each prediction
into per-feature contributions that sum exactly to the result, so it
reports the model's actual arithmetic rather than an approximation of
its reasoning.

\textbf{Day~+1 (XGBoost).} Top drivers are PM\textsubscript{2.5} (3.44),
the six-hour trend (1.83), the 24-hour temperature range (1.81), and
recent AQI lags. The presence of the temperature range is a useful
confirmation --- it was added on physical reasoning about inversions,
and the model independently found it informative.

\textbf{Days~+2 and~+3 (Ridge).} Dominated by \texttt{aqi\_lag\_1h}
(26.32 at day~+2), an order of magnitude above anything else. This is
the persistence effect made visible, and it corroborates the modest gain
over the persistence baseline at those horizons.

The contrast between the two model types is itself informative: the tree
model distributes attention across pollutants and meteorology, while the
linear models concentrate almost entirely on the most recent reading.

Two explanation types are produced --- global importance averaged over
recent forecasts, and a local decomposition of the current prediction
--- and both are surfaced in the dashboard's Accuracy view.

% ---------------------------------------------------------------
\section{Automation and reliability}

\subsection{Scheduled pipelines}

Two GitHub Actions workflows. The hourly one collects new readings,
builds features and writes to both the local CSV and the feature store.
The nightly one collects, retrains all models, registers them, rebuilds
the dashboard snapshot, and commits it back to the repository.

The collector fetches a \emph{range} rather than a single hour: it asks
each destination for its newest timestamp and retrieves everything
since. This is one code path that handles the normal case (one hour) and
recovery after an outage (several hundred) identically, and it means a
missed run needs no special handling.

\subsection{Observed platform limitations}

Three were encountered and worked around; all are documented because
they shaped the design.

\textbf{GitHub scheduling is best-effort.} On free accounts, hourly
crons fire perhaps three or four times a day rather than 24. This is
widely reported and is not specific to this project. The range-fetch
design absorbs it --- data completeness is unaffected, only latency.

\textbf{Hopsworks free-tier reads are unreliable.} Row reads
intermittently time out after roughly 170~seconds. Every read is
therefore retried with growing pauses, and every consumer can fall back
to a local CSV.

\textbf{Large inserts fail materialisation.} Uploading several hundred
rows at once caused the Spark job that makes rows queryable to fail,
leaving them accepted but invisible. Chunking to 100 rows with pauses
resolved it, and the chunk size is now configured accordingly.

\subsection{Degradation behaviour}

The system was designed so that failures degrade rather than break. The
behaviour differs between a developer machine and a scheduled run,
because the two have different local state, and both cases are set out
below.

\begin{center}
\small
\begin{tabular}{lll}
\toprule
\textbf{Failure} & \textbf{On a developer machine} & \textbf{On a scheduled runner} \\
\midrule
Store read fails, collector
  & Uses the local CSV for context
  & Exits; next run retries \\
Store read fails, training
  & Falls back to the local CSV
  & Refuses to train (see below) \\
Upload fails
  & Data already written locally
  & Data already written locally \\
Materialisation deferred
  & Rows accepted; visible later
  & Rows accepted; visible later \\
Scheduled run skipped
  & ---
  & Next run collects the backlog \\
Dashboard cannot reach store
  & Serves the committed snapshot
  & Serves the committed snapshot \\
\bottomrule
\end{tabular}
\end{center}

The second row deserves comment. A developer machine holds the full
48,710-row feature file, so falling back to it is a genuine recovery. A
scheduled runner starts with no local data, and although the collector
writes a file during the run, it holds only the hours just fetched ---
falling back to that would mean training on a few dozen rows.

Training therefore refuses to run on fewer than 5,000 rows. The guard
was added after a scheduled run trained on 42: the runner's freshly
written file carried a newer timestamp than the store and won a
comparison that considered recency alone. Selection now weighs recency
\emph{and} completeness, with the row-count floor as an independent
second check. When it triggers, the run fails visibly and the registered
models stay in service --- the correct outcome, since a stale model
beats one trained on an unrepresentative sample.

The dashboard is unaffected by any of this: it reads a committed
snapshot rather than querying the store, so a failure upstream leaves it
serving the last good forecast, with the issue time shown in the footer.

% ---------------------------------------------------------------
\section{The dashboard}

Deployed on Streamlit Community Cloud with four views: current
conditions and a 72-hour overview; a forecast view offering either a
full day or any single hour; model accuracy with SHAP explanations; and
30 days of history.

The dashboard reads a committed JSON snapshot rather than querying the
feature store live. This was a deliberate reversal of an earlier design.
A live-fetch path existed at first, but it could not produce anything
newer --- the snapshot is rebuilt from the same models in the same
nightly run, so a live read returns the same forecast several seconds
slower, and during a Hopsworks outage it froze the page for minutes.
Reading a committed file cannot fail.

Hazardous-AQI alerting is implemented as a banner that appears whenever
any forecast hour reaches AQI 150 or above, naming the peak value, its
category, and when it occurs.

% ---------------------------------------------------------------
\section{Limitations and further work}

\textbf{Day~+3 accuracy.} At three days the mean absolute error reaches
26 AQI points, wide enough to cross a category boundary. Seven distinct
approaches were tried and all converged on a similar figure, which
suggests the constraint is the information available rather than the
method applied.

\textbf{Observed rather than forecast weather.} The models use the
current day's weather to predict future AQI, where forecast weather for
the target day would be correct. The archive available at this location
was too sparse to use. A paid forecast archive would likely improve
days~+2 and~+3 most.

\textbf{No upwind data.} Ghotki's air is partly transported from
elsewhere. Pollutant readings from Sukkur, Larkana or Rahim Yar Khan,
combined with wind direction, would add genuinely new information rather
than rearranging what is already present. This is the most promising
remaining avenue.

\textbf{Sparse coverage of extreme events.} Only 1.6\% of hours fall in
the Hazardous category, so the models see few examples of exactly the
conditions where a forecast matters most.

\textbf{Dropped-row bias.} Around 4\% of rows are excluded for
incomplete features. Their mean AQI is 110 against 132 for retained rows
--- the exclusions skew towards cleaner periods, so reported accuracy is
if anything slightly conservative.

% ---------------------------------------------------------------
\enlargethispage{6\baselineskip}
\section{Conclusion}

A complete forecasting system was built and deployed, meeting the
specified architecture at every stage and running unattended: hourly
collection, nightly retraining, versioned model registration, and a
served dashboard.

Two decisions shaped the outcome more than any choice of algorithm. The
first was treating the target variable as something to be validated in
its own right. The AQI is computed rather than measured, and a computed
target places a ceiling on every model trained against it --- one no
amount of model comparison will reveal, because every model sits below
it equally. Verifying the implementation against the EPA standard
established that ceiling before any modelling began.

The second was reporting every result against a persistence baseline.
Air quality is strongly autocorrelated, so a model can post a high
R\textsuperscript{2} while contributing little of its own. The gap
between model and baseline is the honest measure, and it widens from
$+0.024$ at one day to $+0.069$ at three --- the models earn more of
their keep precisely where forecasting is hardest. Both practices cost
far less than the modelling work they informed.

\end{document}
